\documentclass[journal]{IEEEtran}

\usepackage{amsmath,amssymb,amsfonts}
\usepackage{graphicx}
\usepackage{booktabs}
\usepackage{multirow}
\usepackage{algorithm}
\usepackage{algorithmic}
\usepackage{hyperref}
\usepackage[dvipsnames,table]{xcolor}
\usepackage{cite}
\usepackage{tabularx}
\usepackage{subcaption}
\usepackage{array}
\usepackage{makecell}
\usepackage{tikz}
\usetikzlibrary{positioning, arrows.meta, decorations.pathreplacing, calc, fit, shapes.geometric}
\usepackage{fontawesome5}

\input{math_commands}

\begin{document}

\title{SpecAnchor: Spectral Anchor Routing for Cross-Domain Class-Incremental Learning with Hyperspectral and LiDAR Data}

\renewcommand{\methodname}{SpecAnchor}
\renewcommand{\backbonename}{SpectralAnchorNet}

\author{Junchao Yang}

\maketitle

\begin{abstract}
Deploying land-cover classification systems across geographically diverse regions demands models that incrementally learn new classes from new domains without forgetting and with minimal additional storage overhead.
We study this \emph{cross-domain multimodal class-incremental learning} (CIL) problem on HSI+LiDAR data spanning three geographic domains (MUUFL, Trento, Houston~2013), 9 sequential tasks, and 32 classes.
Through systematic drift analysis with six architectures of varying branch coupling, we reveal a previously unreported \emph{architecture-dependent drift asymmetry}: in independent-branch fusion backbones, spatial features drift up to 5.8$\times$ more than spectral features across domains, whereas this asymmetry vanishes in tightly coupled architectures.
We identify three correlated factors---branch independence, parameter capacity asymmetry, and feature coupling---and argue that this asymmetry constitutes a \emph{structural prior} of independent-branch designs that can be exploited as a design signal for non-uniform adaptation.
Based on this insight, we propose \methodname{} (Cross-Modal Compensated Drift LoRA), which adopts a tri-branch Mamba backbone with a separate spectral stream and Siamese spatial processing to preserve the exploitable asymmetry, freezes the spectral branch as a stable anchor after warm-up, and applies asymmetric low-rank adapters to the two spatial branches with rank matched to each branch's effective adaptation dimensionality.
A per-domain whitening normalization (SHINE) and drift-compensated prototype reconstruction (DCR) further address cross-domain feature misalignment.
On the 9-task benchmark (32 classes across 3 domains: 11+6+15), \methodname{} achieves $84.5\% \pm 0.9$ balanced accuracy \emph{without maintaining a replay buffer}, storing only compact prototypes, adapter weights, and normalization statistics (${\sim}75\times$ less extra storage than replay buffers) while leveraging within-domain training data for adapter fine-tuning and prototype recomputation---a natural assumption in fixed-sensor remote sensing deployments. See Section~V for the full data access protocol.
It surpasses all replay-free baselines by over $40$\,pp and is competitive with the best replay methods (iCaRL $77.9\%$, FOSTER $77.0\%$ on their original backbones).
Evaluation across all six domain orderings confirms consistent advantages over replay baselines ($80.9\%$ vs.\ $77.8\%$ and $77.2\%$, winning 5 of 6 orderings).
\end{abstract}

\section{Introduction}
\label{sec:intro}

\IEEEPARstart{M}{ultimodal} classification with hyperspectral imagery (HSI) and LiDAR data plays an important role in remote sensing applications such as urban planning, forest monitoring, and geological mapping.
HSI provides dense spectral measurements for material discrimination, whereas LiDAR contributes height and structural cues that help distinguish spectrally similar land-cover classes, making the two modalities complementary~\cite{hong2022multimodal}.
Building on this complementarity, recent HSI-LiDAR fusion networks have achieved strong performance in multimodal remote sensing classification with feature-level and decision-level fusion strategies~\cite{hang2020coupled,fang2022s2enet,li2022mahidfnet}.

With the rapid growth of remote sensing data and the increasing diversity of sensing platforms, land-cover observations accumulate over time across new regions, sensors, and acquisition conditions, giving rise to a lifelong learning scenario in which both the label space and the data domain evolve after initial deployment.
Continual learning studies how models can learn from sequentially arriving data while retaining knowledge acquired from previous tasks~\cite{delange2021continualsurvey,wang2024continualsurvey}.
Typical continual learning settings~\cite{vandeven2022three} include class-incremental learning (CIL), where new categories arrive over time and the model predicts over all classes seen so far, typically within a shared data domain~\cite{rebuffi2017icarl,hou2019lucir}; domain-incremental learning (DIL), where the input distribution shifts across tasks while the label space and task structure remain unchanged~\cite{wang2024pina}; and task-incremental learning (TIL), where the stream is partitioned into tasks and the task identity is provided at inference~\cite{serra2018hat}.
This paper focuses on cross-domain class-incremental learning (CDCIL), where class-incremental label growth and domain-incremental distribution shift are coupled within a single task stream, as illustrated in Fig.~\ref{fig:cdcil_setting}.
  
\begin{figure*}[t!]
  \centering
  \includegraphics[width=\textwidth]{figures/fig1/fig1.pdf}
  \caption{\textbf{Cross-domain class-incremental learning setting on HSI and LiDAR data.}
  The model is trained sequentially on $9$ tasks across three sensors (MUUFL, Trento, and Houston 2013), with the class set accumulating across tasks ($4 \to 32$ classes total). Each task introduces a disjoint subset of land-cover classes; the cumulative class set after task $t$ is denoted ``Class set $t$''.}
  \label{fig:cdcil_setting}
\end{figure*}

CDCIL inherits the core difficulties of both CIL and DIL, but their coupling creates new challenges that neither setting addresses in isolation.
As shown in Fig.~\ref{fig:cdcil_setting}, each domain is divided into a sequence of class-incremental tasks, where every task introduces a subset of new land-cover classes and the cumulative label space grows within that domain.
After the tasks of one domain are completed, the stream moves to a different sensor and location; changes in wavelength range, band count, and ground sample distance (GSD) shift the input distribution while class accumulation continues.
The model therefore faces a sharpened stability--plasticity dilemma~\cite{mermillod2013stability}: it needs to avoid catastrophic forgetting of old classes from previously seen domains, absorb new classes, and align feature representations across changing domains simultaneously.
This forgetting risk arises at two levels: intra-domain forgetting, where decision boundaries fitted to new classes reduce the separability of earlier classes from the same domain; and inter-domain forgetting, where adaptation to a new sensor or location shifts these boundaries away from earlier-domain classes.
These boundary shifts are amplified by semantic overlap across domains.
For example, Fig.~\ref{fig:cdcil_setting} shows that the class ``Trees'' appears in both MUUFL and Houston 2013, ``Apple Trees'' appears in Trento, and ``Road'' appears in all three domains.
The feature space therefore has to support discrimination at two levels: the model has to separate domains with different sensing conditions while distinguishing semantically related classes within and across those domains.

Current methods addressing catastrophic forgetting can be broadly grouped into five directions: replay-based~\cite{rebuffi2017icarl,hou2019lucir,douillard2020podnet,wang2022foster}, regularization-based~\cite{kirkpatrick2017ewc,li2017lwf}, architecture-based~\cite{wang2022l2p,wang2022dualprompt,smith2023codaprompt,liang2024inflora}, meta-learning-based~\cite{yang2026pmpt,zhu2025pass}, and feature-space-based~\cite{petit2023fetril,mcdonnell2023ranpac,gomezvilla2024ldc}.
Remote-sensing CIL methods extend these strategies to hyperspectral and multimodal classification scenarios, but are mostly evaluated within a single scene, sensor, or dataset~\cite{deng2024dcprn,xu2025crossacl,zhang2024jfi_acl,li2026dfpem,yang2026pmpt}.
CDCIL has also attracted attention in other fields, such as medical diagnosis and industrial fault diagnosis, where new categories arrive under changing domains~\cite{yang2023cdfscil,shi2024cdcibn}.
In remote sensing, the closest existing settings sit one or two axes away from ours: HyperKD~\cite{li2025hyperkd} performs single-modality HSI lifelong learning with each dataset as one coarse task; CDCIRSC~\cite{zhang2024cdcirssc} targets scene-level cross-domain CIL on different data sources; and recent HSI-LiDAR transfer methods~\cite{liu2025prompt} adapt across datasets through joint rather than incremental training.
Multimodal HSI-LiDAR CDCIL with finer per-domain task granularity therefore exposes representation-level difficulties that remain insufficiently captured by existing settings, especially when class boundaries evolve within each domain while sensor and scene conditions shift across domains.
This leaves open how to preserve separability between old and new classes within each domain, prevent forgetting across domains, and maintain HSI-LiDAR representations that distinguish semantically similar classes while accounting for domain-specific sensing conditions.

To diagnose these representation-level difficulties, we examine branch-level drift in representative HSI-LiDAR fusion backbones, including FusAtNet~\cite{mohla2020fusatnet}, S2ENet~\cite{fang2022s2enet}, Coupled CNNs~\cite{hang2020coupled}, MAHiDFNet~\cite{li2022mahidfnet}, and HCT~\cite{ghosh2023hct}.
We find that, across these backbones, the spatial-to-spectral drift ratio ranges from $0.60\times$ for FusAtNet to $5.66\times$ for HCT.
Coupled fusion architectures such as FusAtNet show weak or no spectral--spatial drift separation, whereas backbones with more explicit spectral--spatial separation, such as HCT, tend to allocate more drift to the spatial branches.
Motivated by this observation, instead of suppressing branch drift uniformly, we design \backbonename{}, a tri-branch backbone that isolates a capacity-bounded spectral stream as a stable anchor across domains while letting the spatial branches absorb domain-specific variation.
Within this pool, \backbonename{} attains a spectral drift of $0.081$ and a $7.05\times$ spatial-to-spectral ratio, supporting the intended anchor behavior.
Building on this anchor, \methodname{} is a lightweight two-layer CDCIL framework that trains the backbone during first-domain warm-up and keeps it frozen; subsequent domains are processed by forward feature extraction, a compact frozen-feature memory, and incremental updates of domain-specific heads.
The spectral branch routes each sample to a domain-specific head, while the spatial branches contribute drift-suppressed features for fine-grained classification within the routed domain.

The main contributions of this paper are summarized as follows:
\begin{itemize}
    \item We propose \methodname{}, a lightweight two-layer framework specifically tailored to multimodal HSI-LiDAR cross-domain class-incremental learning. It trains the backbone during first-domain warm-up and keeps it frozen; subsequent domains are processed by forward feature extraction, a compact frozen-feature memory, and incremental updates of domain-specific heads. Stable spectral features distinguish domains, whereas drift-suppressed spatial features support fine-grained class discrimination, thereby mitigating intra-domain and inter-domain forgetting under a sharpened stability--plasticity dilemma.
    \item \methodname{} is built on \backbonename{}, a tri-branch backbone motivated by a branch-level drift analysis across five representative HSI-LiDAR fusion backbones. The analysis suggests two design choices associated with stronger spectral--spatial drift asymmetry---isolating the spectral stream from cross-modal coupling and bounding its capacity---and \backbonename{} instantiates both, providing the low-drift spectral anchor used by \methodname{} for domain routing.
    \item We conduct comprehensive CDCIL experiments on MUUFL, Trento, and Houston 2013 with $9$ tasks and $32$ classes, covering baseline comparisons, branch-level drift analysis, spectral-routing evaluation, memory-footprint analysis, and ablation studies to verify the roles of the spectral anchor and drift-suppressed spatial features.
\end{itemize}

\section{Related Work}
\label{sec:related}

\subsection{HSI+LiDAR Multimodal Classification}

Deep learning approaches for joint HSI+LiDAR classification have evolved from early CNN-based methods~\cite{chen2016deep3dcnn} to diverse fusion architectures.
Coupled CNNs~\cite{hang2020coupled} use partially shared convolutional streams for HSI and LiDAR.
FusAtNet~\cite{mohla2020fusatnet} applies cross-modal attention with a shared feature hierarchy.
S2ENet~\cite{fang2022s2enet} employs cross-enhancement modules that tightly couple spectral and spatial features.
MAHiDFNet~\cite{li2022mahidfnet} maintains independent 1D spectral and 2D spatial CNN branches.
HCT~\cite{ghosh2023hct} combines 3D convolutions with transformer fusion in a two-stream architecture.
More recently, state-space models (Mamba)~\cite{gu2024mamba} have shown promise for efficient long-sequence spectral modeling.
Our \backbonename{} backbone adopts a tri-branch design---1D Mamba for spectral, Siamese 2D VSSBlocks for HSI-spatial and LiDAR-spatial---maintaining a separate spectral stream to preserve spectral--spatial drift independence.
None of these fusion architectures consider incremental class arrival or cross-domain deployment.

\subsection{Class-Incremental Learning}

CIL methods fall into three families:
(i)~\emph{regularization-based} methods such as EWC~\cite{kirkpatrick2017ewc} and LwF~\cite{li2017lwf} that constrain weight updates to preserve old-task knowledge;
(ii)~\emph{replay-based} methods such as iCaRL~\cite{rebuffi2017icarl}, LUCIR~\cite{hou2019lucir}, PODNet~\cite{douillard2020podnet}, and FOSTER~\cite{wang2022foster} that store or generate exemplars; and
(iii)~\emph{analytic/prototype-based} methods such as ACIL~\cite{zhuang2022acil} and PASS++~\cite{zhu2025pass} that avoid replay entirely by maintaining class statistics or generating pseudo-features.

In the HSI domain, DCPRN~\cite{deng2024dcprn} combines dual knowledge distillation with prototype replay, and CrossACL~\cite{xu2025crossacl} applies analytic continual learning with cross-projection.
Most recently, PMPT~\cite{yang2026pmpt} applies prototype-based meta-prompt tuning for few-shot class-incremental learning on multimodal remote sensing data, achieving strong results on single-domain benchmarks with pre-trained ViT backbones.
However, existing HSI-CIL methods, including PMPT, operate on individual datasets rather than across geographic domains with heterogeneous sensors, and assume access to a strong pre-trained backbone---conditions that differ fundamentally from our cross-domain buffer-free setting.

\subsection{Parameter-Efficient Continual Learning}

Parameter-efficient fine-tuning (PEFT) adapts pre-trained models with minimal parameter overhead, making it naturally suited to CIL where catastrophic forgetting scales with the number of updated parameters.
Prompt-based approaches---L2P~\cite{wang2022l2p}, DualPrompt~\cite{wang2022dualprompt}, CODA-Prompt~\cite{smith2023codaprompt}---learn task-specific prompts for a frozen ViT backbone.
LoRA-based methods---InfLoRA~\cite{liang2024inflora}, SD-LoRA~\cite{sdlora2025}---apply low-rank updates with orthogonal or decoupled constraints.
EASE~\cite{zhou2024ease} employs expandable adapters with prototype complement.
These methods assume a single-modality backbone (typically an ImageNet-pretrained ViT) and apply \emph{uniform} adaptation across all layers.
These methods are not directly comparable to our setting without nontrivial architectural redesign: (1)~their ViT patch embeddings are designed for 3-channel RGB input and would require re-engineering for 36-band HSI $+$ 2-channel LiDAR, and (2)~they rely on large-scale pre-training (ImageNet-21k) for the frozen backbone to provide meaningful features---no such pre-trained foundation model currently exists for HSI+LiDAR multimodal data.
In contrast, \methodname{} exploits the multimodal structure of HSI+LiDAR to make \emph{modality-asymmetric} adaptation decisions: the spectral branch is frozen as a stable anchor while the two spatial branches receive low-rank adapters with rank matched to each branch's effective adaptation dimensionality.

\subsection{Feature Normalization for Domain Alignment}

Domain-specific batch normalization (DSBN)~\cite{chang2019dsbn} maintains per-domain running statistics to reduce domain shift.
Feature whitening~\cite{pan2019switchable} decorrelates features to remove style-specific information.
In CIL, post-hoc normalization has received less attention, though recent work on multi-prototype methods~\cite{multiproto2025} shows that feature calibration can substantially improve prototype-based classifiers.
Our SHINE module follows this line: it computes per-domain, per-branch whitening statistics and applies Z-score normalization before cosine prototype matching, effectively eliminating the distributional gap between features from different geographic domains.

% sections/3_observation.tex is intentionally excluded: its content is folded
% into the Method motivation to front-load the proposed framework.
\section{Method: \methodname{}}
\label{sec:method}

We present \methodname{} (Cross-Modal Compensated Drift LoRA), a parameter-efficient buffer-free framework for cross-domain multimodal CIL.
The method is designed around the structural prior identified in Section~\ref{sec:observation}: in independent-branch fusion architectures, the spectral branch is naturally stable across domains while spatial branches require adaptation.
Fig.~\ref{fig:overview} illustrates the overall architecture.

\input{figures/fig_framework}

\subsection{\backbonename{} Backbone}
\label{sec:backbone}

\backbonename{} (Spectral-Spatial Contrastive Mamba) is a tri-branch multimodal backbone that processes HSI+LiDAR input through three feature streams:

\begin{itemize}
    \item \textbf{Spectral branch}: A 1D Mamba selective state-space model~\cite{gu2024mamba} processes per-pixel spectral vectors of the HSI input, capturing material absorption signatures via efficient long-sequence modeling. Output: $\fspec \in \R^{d}$.
    \item \textbf{HSI-spatial branch}: After a modality-specific input projection, HSI patches are processed by a stack of 2D VSSBlocks (depthwise convolution + channel MLP with state-space mixing) that extract spatial texture patterns. Output: $\fhsi \in \R^{d}$.
    \item \textbf{LiDAR-spatial branch}: LiDAR elevation/intensity maps pass through an independent input projection and then the \emph{same} VSSBlock stack in a Siamese configuration (shared spatial parameters, independent inputs). Output: $\flid \in \R^{d}$.
\end{itemize}

A lightweight cross-fusion module exchanges information between the two spatial streams after the final VSSBlock via residual cross-modal projections, before global average pooling produces the branch-level features.
The three outputs are concatenated into $f = [\fspec;\, \fhsi;\, \flid] \in \R^{3d}$ for classification.
The architecture uses $d{=}64$, 4 VSSBlocks per spatial branch, and 2 Mamba layers in the spectral branch.\footnote{All reported results use the Mamba~\cite{gu2024mamba} and VMamba VSSBlock implementations. The released code includes fallback paths to standard Transformer blocks when these dependencies are unavailable; fallback paths were not used for any reported experiment.}

The key architectural property is \emph{spectral--spatial independence}: the spectral branch shares no parameters with the spatial processing pipeline.
Section~\ref{sec:observation} shows that this separation preserves the drift asymmetry that \methodname{} exploits---the spectral branch naturally maintains stable representations while spatial branches absorb domain-specific variation, creating the structural prior that enables asymmetric adaptation.
The two spatial branches share VSSBlock weights (Siamese), but after warm-up the backbone is frozen and each branch receives its own independent LoRA adapter (Section~\ref{sec:lora}), enabling modality-specific adaptation despite the shared backbone.

\subsection{Phase 1: Warm-Up Training}
\label{sec:warmup}

During the first $W$ tasks (covering the initial geographic domain), the \backbonename{} backbone is fully trained to establish robust feature representations.
Training uses cross-entropy loss computed from cosine-similarity logits against per-class prototypes, with a reduced learning rate for the spectral branch ($0.1\times$ the base rate) to encourage early stabilization of spectral representations.
After the warm-up phase, the \emph{entire} backbone is frozen---including all three branches.
The spectral branch, having learned domain-invariant absorption signatures, becomes the permanent stable anchor.

The warm-up length $W$ is set to the number of tasks in the first domain (e.g., $W{=}3$ for MUUFL with 3 sub-tasks), ensuring complete initial-domain coverage before freezing.

\subsection{Asymmetric LoRA Adaptation}
\label{sec:lora}

For each task $t \geq W$, the frozen backbone is adapted through low-rank updates applied \emph{only to the two spatial branches}, leaving the spectral anchor untouched.

\paragraph{LoRA insertion.}
After each VSSBlock in the spatial branches, a low-rank output adapter modifies the block's feature map:
\begin{equation}
    h' = h + \alpha \cdot W_{\mathrm{up}} W_{\mathrm{down}} h
\label{eq:lora}
\end{equation}
where $h$ is the VSSBlock output, $W_{\mathrm{down}} \in \R^{r \times d}$ projects to a low-rank bottleneck, $W_{\mathrm{up}} \in \R^{d \times r}$ projects back (implemented as $1{\times}1$ convolutions for spatial feature maps), $r$ is the rank, and $\alpha{=}1.0$ is a fixed scaling factor.
Although the two spatial branches share frozen VSSBlock weights, each branch has its own independent set of LoRA adapters, enabling modality-specific adaptation.

\paragraph{Asymmetric rank allocation.}
The two spatial branches receive different LoRA ranks:
\begin{equation}
    r_{\mathrm{HSI}} = r_0, \quad r_{\mathrm{LiDAR}} = 2 r_0
\label{eq:rank}
\end{equation}
The LiDAR-spatial branch receives $2\times$ the rank because it requires a higher-dimensional adaptation subspace: as the sole carrier of elevation information with no spectral branch backup, LiDAR adaptation spans more diverse directions than HSI-spatial, whose information is partially redundant with the frozen spectral anchor.
This is confirmed by the rank allocation comparison (Table~\ref{tab:rank_allocation}), which shows that our asymmetric allocation outperforms symmetric allocation at the same total rank budget.

\paragraph{Domain-conditioned reuse.}
Rather than allocating a separate adapter per task, we share adapters across tasks within the same geographic domain.
Specifically, tasks from the same domain share a single adapter set.
For the 9-task benchmark (3~domains, 3~tasks each), this yields 2 adapter sets for 6 post-warmup tasks, reducing parameter count by $3\times$ over per-task allocation while acting as an implicit domain-level regularizer.

\paragraph{Domain-selective inference.}
At test time, only the adapter corresponding to the query's domain is activated:
\begin{equation}
    h' = h + \alpha\, W_{\mathrm{up}}^{(d)} W_{\mathrm{down}}^{(d)}\, h
\label{eq:domain_selective}
\end{equation}
where $d$ is the domain index of the test sample and $h$ is the frozen backbone feature.
Each domain's adapter is calibrated for its own distributional characteristics; activating only the relevant adapter avoids interference from adapters trained on unrelated domains.
During \emph{training}, when learning a new domain $k$, a fresh adapter is allocated and fine-tuned on the new domain's data; when learning additional tasks within the same domain, the existing domain adapter is unfrozen and continues training.
During \emph{evaluation}, the domain identity is known because test sets are physically separate per geographic region (see SHINE discussion in Section~\ref{sec:shine}).

\subsection{SHINE: Per-Domain Whitening Normalization}
\label{sec:shine}

Cross-domain prototype matching suffers from distributional misalignment: prototypes computed from one domain's training features may not be directly comparable to query features from another domain due to differences in feature statistics (mean, variance).
SHINE addresses this by applying per-domain, per-branch whitening normalization at both prototype construction and inference time.

\paragraph{Construction.}
After training on each task, for each branch $b$ and each domain $d$, we compute the running mean $\mu_d^b$ and standard deviation $\sigma_d^b$ from the training features of that domain.
Class prototypes are stored \emph{after} Z-score normalization:
\begin{equation}
    \bar{f}_b = \frac{f_b - \mu_d^b}{\sigma_d^b + \epsilon}
\label{eq:shine_norm}
\end{equation}
where $d$ is the domain from which the features originate.
When the sample count for a domain--branch pair is below a threshold ($n < 100$), the normalized features are blended with the original features using $\alpha = n/100$ to avoid unreliable statistics from small samples.

\paragraph{Inference.}
In our benchmark experiments, we use a \emph{domain-aware} mode: the query's domain identity is known because test sets are physically separate per domain, and the query is normalized with its domain's statistics.
In open deployment where domain labels are unavailable, a \emph{domain-agnostic} extension could normalize the query using each candidate domain's statistics in turn and select the highest-scoring class across all domain--prototype pairs.
We leave validation of this domain-agnostic mode to future work; the current evaluation uses domain-aware normalization throughout.

\paragraph{Deployment protocol.}
At training time, per-domain statistics $(\mu_d^b, \sigma_d^b)$ are computed from training features and stored alongside class prototypes (4.5\,KB for 3 domains $\times$ 3 branches $\times$ mean/std $\times$ $d{=}64$).
This post-hoc normalization eliminates the first-order distributional gap (mean and variance shift) between domains without introducing any learnable parameters.
SHINE is an integral component of the \methodname{} framework (analogous to herding selection in iCaRL) and is not applied to baselines.

\subsection{DCR: Drift-Compensated Reconstruction}
\label{sec:dcr}

When spatial adapters are trained on a new domain, the spatial features of old-class prototypes become stale: they were computed with the previous adapter state, which no longer matches the current accumulated adapter configuration.
DCR ensures that old-class prototypes remain aligned with the current accumulated adapter configuration.

\paragraph{Principle.}
After training the adapter for a new domain, old-class prototypes from the \emph{same} domain must reflect the updated adapter state.
For each spatial-branch prototype $\mu_c^b$ ($b \in \{\mathrm{hsi}, \mathrm{lid}\}$) of a class $c$ belonging to domain $d$, the corrected prototype under domain-selective inference is:
\begin{equation}
    \hat{\mu}_c^b = (I + \alpha\, W_{\mathrm{up}}^{(d)} W_{\mathrm{down}}^{(d)})\, \mu_c^b
\label{eq:dcr}
\end{equation}
This ensures prototypes and test features are aligned under the same domain adapter.
Spectral-branch prototypes require no correction since the spectral branch is frozen.
Lemma~1 proves this correction is exact when LoRA acts as a linear additive perturbation at the feature level; in practice, LoRA is inserted within multi-block branches where inter-block nonlinearities introduce a small approximation gap.

\paragraph{Implementation.}
In our experiments, we implement DCR implicitly: after each task, features for all seen classes are re-extracted through the current model (with the domain-specific adapter for each class's domain), and prototypes are recomputed from these updated features.
For a strict deployment scenario without old-data access, Eq.~\eqref{eq:dcr} provides an approximate correction that can be applied directly to stored prototypes; the approximation quality depends on the magnitude of the LoRA update relative to the branch's nonlinearities.

\paragraph{Ablation.}
The ``w/o DCR pipeline'' ablation (Table~\ref{tab:ablation}) disables the full domain-conditioned pipeline: adapter reuse, domain-aware distillation, and prototype correction.
Because these components are architecturally coupled, we report the combined effect ($-1.6$\,pp); see Section~\ref{sec:ablation} for discussion.

\subsection{Multi-Centroid Prototype Classification}
\label{sec:classifier}

Each class is represented by $K{=}2$ centroids per branch, obtained by $k$-means clustering on the training features after each task.
For a test sample $x$ with SHINE-normalized branch features $\{\bar{f}_b\}$, the per-branch score aggregates cosine similarities to the $K$ centroids via log-sum-exp:
\begin{equation}
    s_b(c \mid x) = \log \sum_{k=1}^{K} \pi_{c,k}^b \cdot \exp\!\left(\tau \cdot \cosim(\bar{f}_b,\; \bar{\mu}_{c,k}^b)\right)
\label{eq:multi_centroid_score}
\end{equation}
where $\pi_{c,k}^b = n_{c,k}^b / n_c^b$ is the proportion of class-$c$ training samples assigned to centroid $k$ in branch $b$, $\bar{\mu}_{c,k}^b$ is the SHINE-normalized centroid, and $\tau$ is a temperature parameter.
The final class score aggregates across branches:
\begin{equation}
    s(c \mid x) = \sum_{b \in \{\mathrm{spec}, \mathrm{hsi}, \mathrm{lid}\}} s_b(c \mid x)
\label{eq:ncm}
\end{equation}
The predicted class is $\hat{y} = \argmax_{c \in \mathcal{C}_{1:t}} s(c \mid x)$.

The multi-centroid representation captures intra-class heterogeneity (e.g., a land-cover class spanning two spectral sub-populations), improving discrimination over a single-centroid nearest-class-mean.
This prototype-based classifier avoids the bias toward recent classes that plagues fully-connected heads in CIL settings and requires no additional learnable parameters beyond the stored centroids and mixing weights.

\subsection{GMM-Weighted Orthogonal Projection}
\label{sec:gmm_ortho}

A central challenge in exemplar-free CIL is training LoRA adapters on new-class data \emph{without} degrading old-class features.
When the adapter for domain $d$ is trained on task $t$'s new classes, the low-rank update $\Delta W_b = \alpha\, W_{\mathrm{up}}^{(b)} W_{\mathrm{down}}^{(b)}$ modifies the spatial feature space.
Without constraint, this update can displace old-class prototypes, causing catastrophic forgetting (a na\"ive exemplar-free variant degrades to $33.9\%$; see Section~\ref{sec:limitations}).

We propose a \emph{GMM-weighted orthogonal projection} loss that constrains LoRA updates to preserve old-class prototype directions, using only stored GMM statistics---no old-class data access required.

\paragraph{Formulation.}
For each old class $c$ in the current domain with GMM centroids $\{\mu_{c,k}^b\}_{k=1}^{K}$ and per-dimension standard deviations $\{\sigma_{c,k}^b\}_{k=1}^{K}$ in spatial branch $b \in \{\mathrm{hsi}, \mathrm{lid}\}$, the LoRA perturbation applied to centroid $\mu_{c,k}^b$ is:
\begin{equation}
    \delta_{c,k}^b = \alpha\, W_{\mathrm{up}}^{(b)} W_{\mathrm{down}}^{(b)}\, \mu_{c,k}^b
\label{eq:gmm_perturbation}
\end{equation}
This perturbation should be small, especially along dimensions where the class distribution is tight (low variance = high discriminability).
The importance weight for dimension $j$ is inversely proportional to the variance:
\begin{equation}
    w_{c,k,j}^b = \frac{1}{(\sigma_{c,k,j}^b)^2 + \epsilon}
\label{eq:importance_weight}
\end{equation}
The GMM orthogonal projection loss penalizes the importance-weighted perturbation across all old-class centroids:
\begin{equation}
    \mathcal{L}_{\mathrm{GMM}} = \frac{1}{|\mathcal{S}|} \sum_{c \in \mathcal{C}_{\mathrm{old}}^d} \sum_{k=1}^{K} \sum_{b} \sum_{j=1}^{d} w_{c,k,j}^b \cdot (\delta_{c,k,j}^b)^2
\label{eq:gmm_loss}
\end{equation}
where $\mathcal{C}_{\mathrm{old}}^d$ denotes old classes in the current domain $d$, and $|\mathcal{S}|$ is the total number of centroid--branch terms for normalization.

\paragraph{Geometric interpretation.}
Each GMM centroid defines a direction in feature space that is important for classifying that class.
The variance-weighted penalty ensures that LoRA updates are suppressed along \emph{discriminative} dimensions (where $\sigma^2$ is small and the class is tightly clustered) while allowing adaptation along \emph{non-discriminative} dimensions (where $\sigma^2$ is large and variation is tolerable).
With $K{=}2$ centroids per class and $|\mathcal{C}_{\mathrm{old}}^d|$ old classes in the current domain, the protected subspace spans at most $2 \times |\mathcal{C}_{\mathrm{old}}^d|$ directions per branch in the $d$-dimensional feature space, leaving ample room for new-class adaptation.

\paragraph{Why same-domain only.}
Cross-domain old classes are already protected by a separate mechanism: domain-selective inference (Section~\ref{sec:lora}) ensures that each domain's test features pass through only that domain's adapter.
Combined with frozen backbone BatchNorm (Section~\ref{sec:bn_freeze}), cross-domain features are \emph{exactly invariant} to new-domain LoRA training.
Therefore, $\mathcal{L}_{\mathrm{GMM}}$ only needs to constrain same-domain old-class directions, keeping the protected subspace compact.

\paragraph{Connection to drift asymmetry.}
The GMM orthogonal projection directly exploits the drift asymmetry:
the spectral branch is frozen and requires no protection (its prototypes never change), so $\mathcal{L}_{\mathrm{GMM}}$ operates only on the spatial branches ($b \in \{\mathrm{hsi}, \mathrm{lid}\}$).
The asymmetric rank allocation (Section~\ref{sec:lora}) means the LiDAR branch has more adaptation capacity ($r_{\mathrm{lid}} = 2r_0$), and correspondingly more protected directions from old-class centroids.

\subsection{Exemplar-Free Training Protocol}
\label{sec:ef_protocol}

Combining the above components, \methodname{} achieves true exemplar-free training: each task uses \emph{only its own new-class data}, with no access to any previous task's training samples.

\paragraph{Backbone BatchNorm freezing.}
\label{sec:bn_freeze}
During LoRA training, the backbone is frozen (\texttt{requires\_grad=False}), but a subtle issue arises: calling \texttt{model.train()} switches BatchNorm layers to training mode, causing their running statistics to be updated by the current domain's data.
This pollutes old-domain features at inference time, as BatchNorm running means and variances no longer reflect the old domain's distribution.
We address this by keeping the backbone in evaluation mode during LoRA training:
\begin{equation}
    \texttt{model.train();\quad model.backbone.eval()}
\label{eq:bn_freeze}
\end{equation}
This ensures BatchNorm running statistics remain frozen, making cross-domain feature extraction \emph{deterministic} regardless of which domain was last trained.
We verified this empirically: after the fix, non-current-domain feature statistics are \emph{exactly unchanged} across all tasks (confirmed via A/B comparison).

\paragraph{New-class-only LoRA training.}
For each post-warmup task $t$ in domain $d$, the LoRA adapter is trained using only the new classes introduced by task $t$:
\begin{equation}
    \mathcal{D}_{\mathrm{train}}^{(t)} = \{(x_i, y_i) : y_i \in \mathcal{C}_t\}
\end{equation}
where $\mathcal{C}_t$ is the set of new classes in task $t$.
Old-class preservation is handled entirely through the GMM orthogonal projection (Eq.~\ref{eq:gmm_loss}) and domain-aware feature distillation ($\mathcal{L}_{\mathrm{DKD}}$), neither of which requires old-class data.

\paragraph{Current-domain-only prototype recomputation.}
After training, prototypes and SHINE statistics are recomputed only for the \emph{current domain}.
Cross-domain prototypes remain cached from their last computation and are still valid because:
(i)~the backbone is frozen,
(ii)~BatchNorm running statistics are frozen (Eq.~\ref{eq:bn_freeze}), and
(iii)~domain-selective inference ensures cross-domain features pass through unchanged adapters.
This eliminates all cross-domain data access during the training pipeline.

\paragraph{Data access summary.}
Table~\ref{tab:data_access} summarizes the data access pattern.
At every stage, \methodname{} accesses only the \emph{current task's new-class training data}---the same data that every CIL method must see.
No replay buffer is maintained, and no old-task data is accessed at any point.

\begin{table}[t]
\centering
\caption{Data access per task. \methodname{} only accesses current-task new-class data at every stage.}
\label{tab:data_access}
\small
\begin{tabular}{lll}
\toprule
\textbf{Stage} & \textbf{Data accessed} & \textbf{Old data?} \\
\midrule
LoRA training & Current task new classes & No \\
GMM ortho loss & Stored GMM statistics & No \\
DKD loss & Teacher model snapshot & No \\
Prototype update & Current task new classes & No \\
SHINE update & Current task new classes & No \\
Cross-domain & Cached (unchanged) & No \\
\bottomrule
\end{tabular}
\end{table}

\subsection{Overall Training and Inference}
\label{sec:overall}

\paragraph{Training objective.}
The total loss for task $t \geq W$ combines cross-entropy on new-class data with two regularization terms that protect old classes without requiring old-class data:
\begin{equation}
    \mathcal{L} = \Lce + \lambda_{\mathrm{GMM}} \cdot \mathcal{L}_{\mathrm{GMM}} + \lambda_{\mathrm{DKD}} \cdot \mathcal{L}_{\mathrm{DKD}}
\label{eq:total_loss}
\end{equation}
where $\Lce$ is cross-entropy on the current task's new classes, $\mathcal{L}_{\mathrm{GMM}}$ is the GMM-weighted orthogonal projection loss (Eq.~\ref{eq:gmm_loss}) that constrains LoRA updates to preserve old-class prototype directions, and $\mathcal{L}_{\mathrm{DKD}}$ is a domain-aware feature distillation loss from a teacher snapshot of the previous domain.
DCR prototype correction (Eq.~\ref{eq:dcr}) is applied \emph{after} training for classes whose data was accessed during the current task.

\paragraph{Inference pipeline.}
Given a test sample $x$ from domain $d$:
\begin{enumerate}
    \item Extract $\fspec$, $\fhsi$, $\flid$ from the frozen backbone (BatchNorm in eval mode).
    \item Apply domain-selective LoRA: activate only the adapters associated with domain $d$.
    \item Apply SHINE normalization (Eq.~\ref{eq:shine_norm}) using domain $d$'s statistics.
    \item Compute multi-centroid scores (Eq.~\ref{eq:ncm}) against all stored prototypes.
    \item Predict $\hat{y} = \argmax_c\, s(c \mid x)$.
\end{enumerate}

\paragraph{Parameter budget.}
Each domain adapter pair adds ${\sim}6$K trainable parameters (${\sim}0.20\%$ of the frozen \backbonename{} backbone), with multi-centroid prototypes (2 centroids $\times$ 32 classes $\times$ 3 branches $\times$ 64 dims $\approx$ 48\,KB) and SHINE statistics requiring negligible additional storage---orders of magnitude below the cost of storing exemplar images.

\section{Experiments}
\label{sec:experiments}

\subsection{Experimental Setup}
\label{sec:setup}

\paragraph{Datasets.}
We use three widely adopted HSI+LiDAR benchmarks spanning different geographic regions, sensor configurations, and land-cover vocabularies:
\begin{itemize}
    \item \textbf{MUUFL Gulfport} (Southern Mississippi, USA): 72 HSI bands (380--1050\,nm), 2 LiDAR channels (elevation + intensity), 11 classes, $325 \times 220$ pixels at 1\,m GSD.
    \item \textbf{Trento} (Trento, Italy): 63 HSI bands (400--990\,nm), 1 LiDAR channel (DSM), 6 classes, $600 \times 166$ pixels at 1\,m GSD.
    \item \textbf{Houston 2013} (Houston, TX, USA): 144 HSI bands (380--1050\,nm), 1 LiDAR channel (DSM), 15 classes, $349 \times 1905$ pixels at 2.5\,m GSD.
\end{itemize}

\paragraph{Preprocessing.}
HSI bands are resampled to 36 unified channels via spectral grouping and averaging, harmonizing the heterogeneous band counts across datasets.
LiDAR data is zero-padded to 2 channels.
Input patches of size $9{\times}9$ are extracted centered on each labeled pixel, with per-band standardization applied to HSI and separate standardization for LiDAR.
No data augmentation is applied; all methods receive raw patches to ensure a fair comparison.
The temperature parameter in the multi-centroid scorer (Eq.~\ref{eq:multi_centroid_score}) is set to $\tau{=}10$.

\paragraph{Task structure.}
Each domain is divided into 3 incremental tasks by splitting its class set:
MUUFL [4, 4, 3], Trento [2, 2, 2], Houston [5, 5, 5].
The default domain ordering is MUUFL$\to$Trento$\to$Houston (MTH), yielding $T{=}9$ sequential tasks and $C{=}32$ classes.
Class ordering within each domain is randomized per seed.

\paragraph{Evaluation metric.}
We report \emph{final balanced accuracy}: the per-class mean accuracy (balanced accuracy) computed over all 32 classes at the final task.
All results are averaged over 3 random seeds (0, 1, 2) with standard deviations reported.
Task-wise accuracy evolution across all 9 tasks is shown in Fig.~\ref{fig:taskwise}.

\subsection{Baseline Methods}
\label{sec:baselines}

We compare against 10 baseline methods plus two reference bounds, totaling 12 methods.
Each baseline is implemented following its original paper's backbone architecture, optimizer, learning rate schedule, and hyperparameters.
The \emph{only} adaptation made is the input interface (HSI band count + LiDAR channels) to accommodate the HSI+LiDAR data format.
Each method uses its \emph{complete original pipeline}: iCaRL uses NME with herding exemplars, LUCIR uses cosine classification with weight imprinting, FOSTER uses boosting-compression, and our method uses SHINE normalization with multi-centroid prototype classification.
SHINE is designed for the frozen-backbone prototype pipeline (Section~\ref{sec:shine}) and is part of the \methodname{} method, analogous to how herding exemplar selection is part of iCaRL's pipeline.
The ablation (Table~\ref{tab:ablation}) decomposes the contributions: LoRA alone achieves $79.0\%$ and SHINE alone achieves $79.5\%$, while their combination yields $84.5\%$, demonstrating synergy rather than single-component dominance.

Table~\ref{tab:baselines} summarizes the baseline configurations.

\begin{table}[t]
\centering
\caption{Baseline methods. Each uses its original backbone, hyperparameters, and complete pipeline; only the input interface is adapted for HSI+LiDAR.}
\label{tab:baselines}
\small
\begin{tabular}{clllc}
\toprule
\# & Method & Type & Backbone & Exmp. \\
\midrule
1 & Naive fine-tune & Lower & \backbonename{} & 0 \\
2 & Frozen+proto & Lower & \backbonename{} & 0 \\
3 & EWC~\cite{kirkpatrick2017ewc} & EF & ResNet-32 & 0 \\
4 & LwF~\cite{li2017lwf} & EF & ResNet-32 & 0 \\
5 & ACIL~\cite{zhuang2022acil} & EF & ResNet-32 & 0 \\
6 & PASS++~\cite{zhu2025pass} & EF & ResNet-18 & 0 \\
7 & iCaRL~\cite{rebuffi2017icarl} & EB & ResNet-32 & 20 \\
8 & LUCIR~\cite{hou2019lucir} & EB & ResNet-32 & 20 \\
9 & PODNet~\cite{douillard2020podnet} & EB & ResNet-32 & 20 \\
10 & FOSTER~\cite{wang2022foster} & EB & ResNet-32 & 20 \\
11 & Joint training & Upper & \backbonename{} & all \\
\midrule
--- & \textbf{Ours} & \textbf{BF} & \textbf{\backbonename{}+LoRA} & \textbf{0} \\
\bottomrule
\end{tabular}
\vspace{-1mm}
{\footnotesize EF$=$Exemplar-Free, EB$=$Exemplar-Based (20/class), BF$=$Buffer-Free (no exemplar buffer; re-accesses within-domain training data). Replay budget: $K = 20 \times 32 = 640$ total exemplars, class-balanced herding selection.}
\end{table}

\paragraph{Replay protocol.}
All exemplar-based methods use a unified budget of $K{=}20$ exemplars per class ($20 \times 32 = 640$ total).
After each task, exemplars are rebalanced via class-balanced shrinking ($m = K / |\mathcal{C}_{1:t}|$).
Selection follows herding.
Each method uses its original inference head: iCaRL$\to$NME, LUCIR$\to$cosine classifier, PODNet$\to$as specified, FOSTER$\to$default head.

\paragraph{Hyperparameter protocol.}
Task-level parameters (patch size $9{\times}9$, normalization, augmentation, class ordering, seeds) are unified across all methods.
Method-level parameters (optimizer, scheduler, epochs, learning rate, loss weights, classifier) follow each method's original paper.
If a ResNet-family method does not converge, we uniformly reduce the learning rate (e.g., $0.1 \to 0.01$) for that family---no per-method tuning is performed.

\methodname{} hyperparameters: warm-up uses Adam with $\mathrm{lr}{=}10^{-3}$ for 30 epochs (spectral branch $0.1\times$); LoRA training uses Adam with $\mathrm{lr}{=}5{\times}10^{-4}$ for 50 epochs with cosine annealing.
Loss weights: $\lambda_{\mathrm{proto}}{=}1.0$, $\lambda_{\mathrm{ortho}}{=}0.1$, $\lambda_{\mathrm{DKD}}{=}0.5$ (for all post-warmup tasks; the teacher is the previous domain's model snapshot).
LoRA scaling factor $\alpha{=}1.0$; prototype centroids $K{=}2$; multi-centroid temperature $\tau{=}10$.

\subsection{Main Results}
\label{sec:main_results}

Table~\ref{tab:main} presents the full comparison on the 9-task MTH benchmark.

\begin{table}[t]
\centering
\caption{Main results on the cross-domain HSI+LiDAR benchmark (MUUFL$\to$Trento$\to$Houston2013, 9 tasks, 32 classes). All accuracy and forgetting values are balanced accuracy (per-class mean), computed consistently from per-class predictions. All values averaged over 3 seeds.}
\label{tab:main}
\renewcommand{\arraystretch}{1.1}
\setlength{\tabcolsep}{3pt}
\resizebox{\columnwidth}{!}{%
\begin{tabular}{l c c c c c c c c c}
\toprule
\textbf{Method} & \textbf{Mem.} & \textbf{BB} & \textbf{Avg Acc$\uparrow$} & \textbf{BA\textsubscript{M}} & \textbf{BA\textsubscript{T}} & \textbf{BA\textsubscript{H}} & \textbf{Fgt$\downarrow$} & \textbf{Gain} \\
\midrule
\multicolumn{9}{l}{\textit{\color{gray}Upper Bound}} \\
\color{gray}Joint & \color{gray}all & \color{gray}S2CM & \color{gray}86.7{\scriptsize$\pm$0.6} & \color{gray}72.9 & \color{gray}97.2 & \color{gray}92.6 & \color{gray}6.9 & \color{gray}+19.9 \\
\midrule
\multicolumn{9}{l}{\textit{\color{gray}Lower Bounds}} \\
\color{gray}Frozen\textsuperscript{a} & \color{gray}0 & \color{gray}S2CM & \color{gray}57.3{\scriptsize$\pm$2.7} & \color{gray}52.9 & \color{gray}78.6 & \color{gray}52.0 & \color{gray}25.1 & \color{gray}+3.2 \\
\color{gray}Finetune & \color{gray}0 & \color{gray}S2CM & \color{gray}15.0{\scriptsize$\pm$0.1} & \color{gray}0.0 & \color{gray}0.0 & \color{gray}32.1 & \color{gray}82.8 & \color{gray}$-$11.5 \\
\midrule
\multicolumn{9}{l}{\textit{Replay-Free Methods}} \\
EWC    & 0 & R32 & 13.5{\scriptsize$\pm$0.3} & 0.0 & 0.0 & 28.9 & 83.3 & $-$9.7 \\
LwF    & 0 & R32 & 12.3{\scriptsize$\pm$0.5} & 0.0 & 0.0 & 26.3 & 83.7 & $-$5.9 \\
PASS   & 0 & R18 & 11.2{\scriptsize$\pm$0.7} & 0.0 & 0.0 & 23.9 & 60.5 & $-$14.6 \\
ACIL   & 0 & R32 & 39.3{\scriptsize$\pm$5.5} & 31.2 & 35.2 & 46.8 & 33.8 & +0.6 \\
\midrule
\multicolumn{9}{l}{\textit{Replay-Based Methods (20 exemplars/class)}} \\
iCaRL  & 20/c & R32 & \underline{77.9}{\scriptsize$\pm$1.2} & \underline{68.9} & 87.5 & \underline{80.7} & \underline{12.9} & \underline{+7.6} \\
LUCIR  & 20/c & R32 & 75.9{\scriptsize$\pm$2.1} & 67.7 & 85.6 & 78.1 & 15.2 & +8.1 \\
FOSTER & 20/c & R32 & 77.0{\scriptsize$\pm$0.2} & 67.1 & \textbf{93.2} & 77.7 & 14.8 & +2.4 \\
PODNet & 20/c & R32 & 39.7{\scriptsize$\pm$8.8} & 29.4 & 40.8 & 46.8 & 50.2 & $-$30.3 \\
\midrule
\multicolumn{9}{l}{\textit{Ours}} \\
\rowcolor{red!5}
\textbf{CMCD-LoRA} & \textbf{0} & \textbf{S2CM} & \textbf{84.5}{\scriptsize$\pm$0.9} & \textbf{76.0} & \underline{93.3} & \textbf{87.3} & \textbf{8.7} & \textbf{+10.8} \\
\bottomrule
\end{tabular}%
}

\vspace{2pt}
{\footnotesize Avg Acc\,=\,final balanced accuracy (\%, per-class mean over all 32 classes at the last task). BA\textsubscript{M}/BA\textsubscript{T}/BA\textsubscript{H}\,=\,per-domain balanced accuracy at the end of the stream. Fgt\,=\,per-domain forgetting: mean over domains of (peak BA over the full stream $-$ final BA), lower is better. Gain\,=\,T8$-$T2 balanced accuracy change, reflecting net knowledge accumulation. Mem.\,=\,exemplars/class. R32\,=\,ResNet-32, R18\,=\,ResNet-18-CBAM. $\pm$ values are population standard deviations over $N{=}3$ seeds ($\mathrm{ddof}{=}0$). \textbf{Bold}\,=\,best, \underline{underline}\,=\,2nd best (excl.\ bounds). \textsuperscript{a}Frozen: backbone trained on task\,0 only (4 classes), then frozen with cosine NCM; ablation ``Frozen'' (Table~\ref{tab:ablation}) uses the full 3-task warm-up (11 classes).}
\end{table}


\paragraph{Buffer-free accuracy--storage trade-off.}
Following standard CIL evaluation practice, each baseline uses its original backbone and hyperparameters as prescribed in its respective paper; only the input interface is adapted for HSI+LiDAR.
Under this protocol, \methodname{} achieves $84.5\% \pm 0.9$ balanced accuracy without maintaining a replay buffer, while re-accessing within-domain training data for adapter fine-tuning and prototype recomputation.

\paragraph{Backbone capacity discussion.}
Our method uses the \backbonename{} backbone (3.07M params) while baselines use ResNet-32 (0.47M), so the raw accuracy gap partly reflects backbone capacity.
This is a deliberate design choice, not an incidental advantage: the observation study (Section~\ref{sec:observation}) shows that independent-branch architectures like \backbonename{} enable the drift asymmetry that \methodname{} exploits, whereas ResNet-32 (a single-stream architecture) would not exhibit the structural prior that motivates our method.
The backbone selection is thus \emph{part of the contribution}---our method requires a backbone with specific structural properties, not merely a larger one.
We note that the efficiency advantage of \methodname{} lies specifically in rehearsal buffer elimination (${\sim}75\times$ less extra storage, zero exemplars), not in computational cost: Table~\ref{tab:efficiency} shows that \backbonename{} has higher inference latency than ResNet-32.

\paragraph{Data access protocol.}
\methodname{} does not maintain a replay buffer of raw training samples.
After each task, it stores multi-centroid prototypes, SHINE statistics (per-domain mean/variance), and LoRA adapter weights---totaling $\sim$101\,KB vs.\ 7.5\,MB for replay methods.

\methodname{} re-accesses within-domain training data at two points during the training pipeline:
(1)~During \emph{LoRA training}, when the domain-conditioned reuse mechanism (DCR) encounters a returning domain, the LoRA adapter is fine-tuned on \emph{all seen classes from that domain}---not just the new task's classes. This within-domain rehearsal is essential for prototype--adapter alignment (an exemplar-free variant that restricts training to current-task data only degrades to $33.9\%$; see Section~\ref{sec:limitations}).
(2)~After each task, prototypes, centroid sets, and SHINE statistics are \emph{recomputed} by re-extracting features from all seen training data through the current model.

This within-domain data accessibility is a natural property of fixed-sensor remote sensing deployments, where historical observations from the same geographic site remain available.
The key distinction from exemplar-based methods is \emph{storage}: \methodname{} stores no raw samples as part of the CIL pipeline (${\sim}75\times$ less extra storage), accessing the original training data in place rather than maintaining a separate buffer.

\paragraph{Low forgetting.}
\methodname{} achieves the lowest forgetting among non-joint methods ($8.7\%$), lower than iCaRL ($12.9\%$) and FOSTER ($14.8\%$), and drastically lower than catastrophic-forgetting methods ($>60\%$).
The frozen spectral anchor preserves old-domain spectral signatures, and per-domain retention remains stable across domain transitions (Section~\ref{sec:analysis}).

\paragraph{Knowledge accumulation.}
The Gain column ($\Delta = \mathrm{T8} - \mathrm{T2} = +10.8$\,pp) shows that \methodname{} achieves the \emph{largest} knowledge accumulation among non-joint methods, accumulating substantially more than replay baselines (iCaRL $+7.6$, FOSTER $+2.4$).
While replay methods also show positive gain, \methodname{} does so without maintaining a separate exemplar buffer, and PODNet ($-30.3$\,pp) actively degrades.

\paragraph{Cross-domain degradation of single-domain CIL methods.}
EWC, LwF, PASS, and naive fine-tuning all degrade to $\leq 15\%$ balanced accuracy, indicating that methods designed for single-domain CIL are insufficient for the cross-domain multimodal setting.
Even with regularization (EWC) or distillation (LwF), the magnitude of cross-domain distributional shift exceeds what these mechanisms can compensate.

\subsection{Ablation Study}
\label{sec:ablation}

Table~\ref{tab:ablation} presents a systematic ablation of \methodname{}'s components.

\input{figures/ablation_table}

\paragraph{SHINE and LoRA are complementary (\#2, \#5).}
Removing SHINE reduces accuracy by $5.5$\,pp ($84.5\% \to 79.0\%$), while removing LoRA reduces it by $5.0$\,pp ($84.5\% \to 79.5\%$).
Critically, \emph{neither component alone accounts for the full improvement}: removing both yields $70.1\%$ (frozen baseline), while the combined gain is $14.4$\,pp---demonstrating a \emph{synergistic interaction}.
This synergy arises because SHINE addresses the \emph{distributional gap} between domains (mean/variance shift), while LoRA addresses the \emph{representational gap} (spatial feature drift).
Without LoRA, SHINE can only normalize stale frozen features; without SHINE, adapted features from different domains remain misaligned for prototype matching.
Notably, even without SHINE, LoRA alone ($79.0\%$) is competitive with the best replay method ($77.9\%$), confirming that spatial adaptation provides genuine value beyond normalization.

\paragraph{Domain-conditioned reuse + prototype correction (\#3).}
The ``w/o DCR'' ablation disables the full domain-conditioned pipeline: adapter reuse across same-domain tasks, domain-aware feature distillation ($\lambda_\mathrm{DKD}{=}0$), and prototype correction, reducing accuracy by $1.6$\,pp ($82.9\%$ vs.\ $84.5\%$).
These components are architecturally coupled: adapter reuse determines which adapter state prototypes are computed against, and disabling reuse changes the teacher selection for DKD. Specifically, without reuse each task gets an independent adapter, eliminating the domain-level teacher relationship that DKD relies on. Isolating each factor individually would require redesigning the teacher selection logic; we therefore report the combined effect and note that the modest $1.6$\,pp contribution suggests none of the three sub-components dominates.
Note that within-domain data access (Section~\ref{sec:main_results}) is maintained across all ablation configurations; this ablation isolates the algorithmic components of the DCR pipeline, not the data access protocol.

\paragraph{Asymmetric rank allocation matters (\#4).}
Table~\ref{tab:rank_allocation} compares asymmetric ($r_\mathrm{H}{=}4$, $r_\mathrm{L}{=}8$) versus symmetric ($6,6$) allocation at the same total rank budget ($R{=}12$).
Our asymmetric allocation outperforms the symmetric baseline by $1.1$\,pp ($84.5\%$ vs.\ $83.4\%$), confirming that modality-aware rank allocation provides value.
This is consistent with LiDAR being the sole carrier of elevation information with no spectral branch backup, requiring more adaptation directions despite comparable CKA drift magnitude.

\input{figures/rank_allocation_table}

\paragraph{Decoupled spatial adapters are necessary (\#5).}
Sharing a single LoRA adapter between HSI-spatial and LiDAR-spatial branches (at HSI rank $r_0$) tests whether modality-specific adaptation is needed.
Since HSI-spatial and LiDAR-spatial features capture fundamentally different physical properties (texture vs.\ elevation), a shared adapter cannot simultaneously address both drift patterns.

\paragraph{Spatial adaptation is the right target (\#6).}
Applying LoRA only to the spectral branch (with spatial branches frozen) provides a reverse validation of the core hypothesis.
If the drift asymmetry insight is correct, adapting the already-stable spectral branch while leaving the drifting spatial branches frozen should \emph{degrade} performance---confirming that adaptation effort should target the spatial branches.

\subsection{Analysis}
\label{sec:analysis}

\paragraph{Task-wise accuracy evolution.}
Fig.~\ref{fig:taskwise} shows the Avg TAg progression across all 9 tasks.

\begin{figure}[t]
    \centering
    \includegraphics[width=\columnwidth]{figures/fig_taskwise.pdf}
    \caption{Task-wise accuracy evolution on the MTH benchmark. \methodname{}'s advantage over other replay-free methods widens as new domains are introduced at tasks 3 and 6, confirming that the drift-aware adaptation mechanism is most valuable precisely when cross-domain forgetting is most severe.}
    \label{fig:taskwise}
\end{figure}

As Fig.~\ref{fig:taskwise} shows, \methodname{}'s advantage over replay-free baselines widens as new domains are introduced at tasks 3 and 6, confirming that drift-aware spatial adaptation is most valuable at domain transitions where cross-domain forgetting is most severe.

\paragraph{Classification maps.}
Fig.~\ref{fig:classmaps} shows qualitative classification results across the three domains.

\begin{figure}[t]
    \centering
    \includegraphics[width=\columnwidth]{figures/fig_classmaps.pdf}
    \caption{Classification maps on the three domains after the final task (task 8). Each column shows one domain; rows correspond to ground truth and predictions from selected methods. \methodname{} produces more spatially coherent predictions, particularly in the Houston domain where cross-domain forgetting is most severe.}
    \label{fig:classmaps}
\end{figure}

\paragraph{Per-domain retention analysis.}
Table~\ref{tab:main} reveals domain-specific retention patterns.
On MUUFL (the first-learned domain), \methodname{} retains $76.0\%$ balanced accuracy after learning two subsequent domains, substantially above iCaRL ($68.9\%$) and FOSTER ($67.1\%$).
On Houston (the last domain), \methodname{} achieves $87.3\%$---$6.6$\,pp above iCaRL ($80.7\%$)---demonstrating superior plasticity for new domains.
This retention pattern is a direct consequence of the frozen spectral anchor: the anchor preserves old-domain discriminative spectral signatures while spatial LoRA adapters efficiently acquire new-domain spatial features.

\paragraph{Domain tracking curves.}
Fig.~\ref{fig:taskwise} (3 subplots by domain) shows how each domain's balanced accuracy evolves across all 9 tasks.
\methodname{}'s MUUFL accuracy remains stable across domain transitions, while baselines show progressive degradation after domain switches.

\paragraph{Domain ordering robustness.}
To evaluate sensitivity to the order in which domains are presented, we run all 6 permutations of the 3 domains, each with 3 seeds (Table~\ref{tab:domain_order} and Fig.~\ref{fig:6order}).

\input{figures/6order_table}

\begin{figure}[t]
    \centering
    \includegraphics[width=\columnwidth]{figures/fig_6order_robustness.pdf}
    \caption{Domain ordering robustness. \methodname{} achieves competitive accuracy across orderings without maintaining a replay buffer. Dashed lines indicate 6-order means.}
    \label{fig:6order}
\end{figure}

\methodname{} achieves $80.9\%$ averaged across all 6 orderings, outperforming iCaRL ($77.8\%$, $+3.1$\,pp) and FOSTER ($77.2\%$, $+3.7$\,pp) in 5 out of 6 orderings \emph{without maintaining an exemplar buffer}.
Performance is highest when Houston is learned last (MTH: $84.5\%$, TMH: $84.6\%$), where warm-up occurs on the two larger domains (MUUFL: 51K samples, Trento: 29K) that provide richer training data before backbone freezing.
Even in the worst-case ordering (MHT: $78.0\%$), \methodname{} remains at parity with replay baselines, confirming robustness across domain orderings.

\paragraph{Rehearsal storage efficiency.}
\methodname{}'s key deployment advantage is the elimination of exemplar buffer storage.
Each LoRA adapter pair (HSI rank-4 + LiDAR rank-8) adds only $\sim$6K parameters per domain---$0.20\%$ of the frozen backbone.
The total extra storage (prototypes + LoRA adapters + SHINE statistics) is $\sim$101\,KB---${\sim}75\times$ less than the $7.5$\,MB needed to store 640 exemplar patches for replay methods.
Note that the S2CM backbone itself (3.07M parameters) is larger and slower than ResNet-32 (0.47M; see Table~\ref{tab:efficiency} in supplementary); the efficiency advantage is specifically in rehearsal buffer elimination.

\paragraph{Limitations and discussion.}
\label{sec:limitations}
We acknowledge two primary limitations.
\emph{(1)~Within-domain data accessibility.}
\methodname{} requires access to within-domain training data for LoRA fine-tuning and prototype recomputation.
A strictly exemplar-free variant that restricts each task to its own new-class data degrades catastrophically ($33.9\%$ vs.\ $84.5\%$), because LoRA adaptation on new classes alone causes prototype--adapter misalignment for old classes within the same domain.
This confirms that within-domain data re-access is essential, not merely beneficial.
In fixed-sensor remote sensing deployments, this assumption is naturally satisfied.
\emph{(2)~Backbone coupling}---the method's effectiveness is tied to independent-branch architectures. However, this is by design: the observation study (Section~\ref{sec:observation}) provides a principled criterion for backbone selection, and the principled selection of backbone and adaptation strategy is itself a contribution.
Regarding domain ordering, the comprehensive 6-ordering evaluation (Table~\ref{tab:domain_order}) confirms that \methodname{} outperforms replay baselines in 5 of 6 orderings without an exemplar buffer, though with larger variation ($\pm 2.7$\,pp) than baselines ($\pm 0.3$--$0.6$\,pp) due to sensitivity to warm-up domain size.

\section{Conclusion}
\label{sec:conclusion}

We presented \methodname{}, a buffer-free framework for cross-domain multimodal class-incremental learning on HSI+LiDAR data.
Through systematic drift analysis across six fusion architectures, we identified an \emph{architecture-dependent drift asymmetry}---a structural prior of independent-branch designs where spatial features drift significantly more than spectral features.
Exploiting this prior, \methodname{} freezes the spectral branch as a stable anchor and applies asymmetric LoRA adapters to the spatial branches with rank matched to each branch's effective adaptation dimensionality.
Combined with SHINE domain whitening and DCR prototype correction, the framework achieves competitive accuracy ($84.5\%$) with replay methods on a 9-task cross-domain benchmark while requiring ${\sim}75\times$ less extra storage than exemplar buffers.
\methodname{} leverages within-domain training data accessibility---a natural property of fixed-sensor deployments---for adapter fine-tuning and prototype recomputation, eliminating the need for a separate exemplar replay buffer.
A strictly exemplar-free variant degrades catastrophically ($33.9\%$), confirming that within-domain data re-access is architecturally essential for prototype--adapter alignment in multi-block LoRA systems.

\paragraph{Limitations.}
The method's effectiveness relies on backbone designs with spectral--spatial separation---a deliberate choice validated by the observation study, but one that restricts generalization to architectures where spectral and spatial streams are not tightly coupled.
The warm-up length is set heuristically; an automatic convergence criterion would improve robustness.
The 6-ordering evaluation (Table~\ref{tab:domain_order}) confirms consistent advantages over replay baselines (winning 5 of 6 orderings), though with variation of $\pm 2.7$\,pp depending on the initial domain's data volume.
Domain-adaptive warm-up strategies could further reduce this sensitivity.

\paragraph{Reproducibility.}
All experiment result artifacts (per-seed JSONs, per-class CSVs), figure generation scripts, and the complete paper source are included in the supplementary material.
The experiment code, trained checkpoints, and instructions for reproducing results on the three HSI+LiDAR benchmarks will be released upon acceptance.

\paragraph{Future work.}
Promising extensions include: (i)~additional sensor modalities (SAR, multispectral) and longer task sequences; (ii)~domain-adaptive warm-up based on drift convergence criteria; and (iii)~applying the drift asymmetry insight to multi-temporal change detection where temporal modality stability varies.


\appendix[Supplementary Material]

\subsection{Supplementary Ablation}
\label{app:supp_ablation}

Table~\ref{tab:supp_ablation} presents two additional ablation configurations that complement the main ablation study.

\begin{table}[h]
\centering
\caption{Supplementary ablation (9 tasks, MTH, 3-seed mean$\pm$std).}
\label{tab:supp_ablation}
\begin{tabular}{lcc}
\toprule
Configuration & Avg TAg (\%) & $\Delta$ \\
\midrule
Full \methodname{}+SHINE & $84.5 \pm 0.9$ & --- \\
\midrule
w/o LoRA (frozen spatial) & $79.5 \pm 1.6$ & $-5.0$ \\
FC head (replace prototype) & $52.4 \pm 10.0$ & $-32.1$ \\
\bottomrule
\end{tabular}
\end{table}

\paragraph{w/o LoRA.}
Completely removing LoRA adaptation (spatial branches remain frozen after warm-up) isolates the contribution of spatial adaptation itself.
The gap between this configuration and the full method quantifies the value of LoRA-based plasticity on the spatial branches.

\paragraph{FC head.}
Replacing the cosine NCM classifier with a fully-connected (FC) head evaluates the importance of prototype-based classification in the CIL setting.
FC heads are known to exhibit strong recency bias in incremental learning, favoring recent classes over old ones.

\subsection{Efficiency Comparison}
\label{app:efficiency}

\begin{table}[t]
\centering
\caption{Efficiency comparison across CIL methods. ``Extra Storage'' denotes the additional \emph{deployment-time} storage beyond the backbone checkpoint (exemplars, prototypes, adapters, etc.). During training, our method additionally maintains domain-specific teacher snapshots (${\sim}12$\,MB per domain); these are discarded after training. All inference times measured on a single RTX 4090 with batch size 64.}
\label{tab:efficiency}
\renewcommand{\arraystretch}{1.1}
\setlength{\tabcolsep}{4pt}
\resizebox{\columnwidth}{!}{%
\begin{tabular}{l c r r r r r r}
\toprule
\textbf{Method} & \textbf{Backbone} & \textbf{Params} & \textbf{Train} & \textbf{Infer.} & \textbf{GPU Mem.} & \textbf{Extra Storage} & \textbf{Avg Acc} \\
 & & (M) & (min) & (ms) & (MB) & & (\%) \\
\midrule
\multicolumn{8}{l}{\textit{Replay-Based (20 exemplars/class)}} \\
iCaRL  & R32      & 0.47 & 17 & 3.7 & 11.9 & 7.51\,MB  & 77.9 \\
LUCIR  & R32      & 0.47 & 17 & 3.7 & 11.9 & 7.51\,MB  & 75.9 \\
FOSTER & R32      & 0.47 & 27 & 3.7 & 11.9 & 9.32\,MB  & 77.0 \\
PODNet & R32      & 0.47 & 17 & 3.7 & 11.9 & 7.51\,MB  & 39.7 \\
\midrule
\multicolumn{8}{l}{\textit{Replay-Free}} \\
EWC    & R32      & 0.47 & 17 & 3.7 & 11.9 & 0.47\,MB\textsuperscript{$\dagger$} & 13.5 \\
PASS   & R18-CBAM & 11.2 & 12 & 2.5 & 62.8 & ---       & 11.2 \\
ACIL   & R32      & 0.47 & 17 & 3.7 & 11.9 & ---       & 39.3 \\
\midrule
\rowcolor{red!5}
\textbf{Ours} & \textbf{S2CM} & \textbf{3.07} & \textbf{${\sim}$54} & \textbf{34.1} & \textbf{71.9} & \textbf{101\,KB} & \textbf{84.5} \\
\bottomrule
\end{tabular}%
}

\vspace{2pt}
{\footnotesize Params includes backbone + task-specific adapters. Extra Storage: replay methods store 640 raw patches ($38{\times}9{\times}9$, float32); FOSTER additionally stores the old-task model (1.8\,MB); Ours stores multi-centroid prototypes (48\,KB: $2 \times 32 \times 3 \times 64 \times 4$\,B for $K{=}2$ centroids) + mixing weights (0.8\,KB) + 2 LoRA adapter sets (48\,KB: 2 domains $\times$ $[r_H{=}4, r_L{=}8]$ $\times$ 4 VSSBlocks) + SHINE statistics (4.5\,KB: $3 \times 3 \times 2 \times 64 \times 4$\,B). \textsuperscript{$\dagger$}EWC stores a Fisher information matrix of equal size to the model. Infer.\,=\,forward pass time per batch of 64 samples. Ours achieves ${\sim}75\times$ less extra storage than replay methods.}
\end{table}


\subsection{Per-Class Accuracy}
\label{app:perclass}

Detailed per-class accuracy after the final task (task 8) is provided in Table~\ref{tab:perclass} for \methodname{} and representative baselines.

\begin{table*}[t]
\centering
\caption{Per-class accuracy (\%) after task 8 (seed 0 for Joint, iCaRL, and Finetune; seed 1 for Ours, selected as the median-accuracy seed). Classes grouped by domain. Seed-averaged results are reported in Table~\ref{tab:main}.}
\label{tab:perclass}
\scriptsize
\setlength{\tabcolsep}{2.5pt}
\begin{tabular}{llrrrr}
\toprule
Domain & Class & \textbf{Joint} & \textbf{Ours} & \textbf{iCaRL} & \textbf{Finetune} \\
\midrule
MUUFL & Trees & 90.8 & 66.8 & 74.8 & 0.0 \\
 & M.Grass & 63.0 & 60.2 & 69.9 & 0.0 \\
 & MixGnd & 64.9 & 66.7 & 58.6 & 0.0 \\
 & Dirt & 73.8 & 64.8 & 80.1 & 0.0 \\
 & Road & 78.4 & 67.2 & 71.6 & 0.0 \\
 & Water & 61.2 & 52.4 & 47.2 & 0.0 \\
 & BldgSh & 61.1 & 75.9 & 74.1 & 0.0 \\
 & Bldg & 92.3 & 87.3 & 87.9 & 0.0 \\
 & Sidew & 52.5 & 51.1 & 66.1 & 0.0 \\
 & YCurb & 75.2 & 83.2 & 61.6 & 0.0 \\
 & Cloth & 78.7 & 76.0 & 50.8 & 0.0 \\
\midrule
Trento & Apple & 97.3 & 99.5 & 98.3 & 0.0 \\
 & Bldg & 99.1 & 89.9 & 91.8 & 0.0 \\
 & Ground & 100.0 & 51.6 & 78.3 & 0.0 \\
 & Woods & 100.0 & 99.8 & 99.7 & 0.0 \\
 & Viney & 98.8 & 53.9 & 92.5 & 0.0 \\
 & Roads & 95.5 & 87.3 & 82.8 & 0.0 \\
\midrule
Houston & HGrass & 82.8 & 84.8 & 81.1 & 0.0 \\
 & SGrass & 99.6 & 96.6 & 96.9 & 0.0 \\
 & SynGr & 100.0 & 99.4 & 90.7 & 0.0 \\
 & Trees & 98.5 & 98.7 & 96.0 & 0.0 \\
 & Soil & 100.0 & 100.0 & 88.9 & 0.0 \\
 & Water & 98.6 & 91.6 & 98.6 & 0.0 \\
 & Resid & 92.3 & 79.9 & 81.2 & 0.0 \\
 & Comm & 93.4 & 90.9 & 89.0 & 0.0 \\
 & Road & 91.5 & 71.2 & 60.2 & 0.0 \\
 & Hwy & 78.9 & 58.8 & 53.9 & 0.0 \\
 & Rail & 92.3 & 78.7 & 89.5 & 99.3 \\
 & PLot1 & 89.2 & 76.0 & 49.6 & 95.6 \\
 & PLot2 & 91.2 & 84.9 & 75.4 & 85.6 \\
 & Tennis & 100.0 & 98.8 & 97.6 & 100.0 \\
 & Track & 96.2 & 100.0 & 85.6 & 100.0 \\
\bottomrule
\end{tabular}
\end{table*}


\subsection{Theoretical Analysis}
\label{app:theory}

We provide formal analysis of the three core design choices in \methodname{} under \emph{local linearization assumptions}.
These results characterize the behavior of the method in the vicinity of the frozen-backbone operating point and are intended as design intuitions, not global guarantees for the full nonlinear training dynamics.
Throughout, we use $f = [\fspec;\, \fhsi;\, \flid] \in \R^{3d}$ for concatenated features from the three branches, and $\mu_c \in \R^{3d}$ for the class-$c$ prototype.

\subsubsection{Lemma 1: DCR Correction Under Linear LoRA}
\label{app:dcr_exact}

\begin{lemma}[DCR is exact under single-layer linear LoRA]
Let $f_b(x)$ denote the frozen spatial-branch feature for input $x$, and let the domain-specific LoRA adapter produce $\tilde{f}_b(x) = (I + \Delta W_b)\, f_b(x)$ where $\Delta W_b = \alpha\, W_{\mathrm{up}} W_{\mathrm{down}}$ is the adapter for the class's domain.
Then the DCR-corrected prototype $\hat{\mu}_c^b = (I + \Delta W_b)\,\mu_c^b$ equals the mean of the adapted features exactly:
\begin{equation}
    (I + \Delta W_b)\,\mu_c^b = \frac{1}{|S_c|}\sum_{x \in S_c} (I + \Delta W_b)\,f_b(x)
\end{equation}
\end{lemma}
\begin{proof}
$(I + \Delta W_b)$ is a fixed linear operator. The mean commutes with linear maps: $(I + \Delta W_b)\,\mathbb{E}[f_b] = \mathbb{E}[(I + \Delta W_b)\,f_b]$.
\end{proof}

This holds for any adapter and any rank.
It breaks down only if adaptation includes nonlinearities (e.g., adapters with activation functions), which standard LoRA does not.
Under the stated single-layer linear assumption, DCR is exact when updating old-class prototypes.
In multi-block architectures where nonlinearities intervene between LoRA layers, this degrades to a first-order approximation with residual error of $O(\|\Delta W\|^2)$.
Nevertheless, the linear correction provides a principled initialization that replay-based methods, which must approximate old-class distributions from a finite exemplar set, cannot guarantee.

\subsubsection{Proposition 1: Independent Branches Reduce Cross-Domain Interference}
\label{app:interference}

This proposition provides intuition for the core architectural insight from Section~\ref{sec:observation}: why independent branches lead to less catastrophic forgetting than shared-parameter designs.

\begin{proposition}[Cross-branch interference under new-domain adaptation]
\label{prop:interference}
Let $\mathcal{R}_{\mathrm{old}}(\theta) = \mathbb{E}_{(x,y) \sim \mathcal{D}_{\mathrm{old}}}[\ell(f_\theta(x), y)]$ be the old-task risk, and let $\theta^*$ be the parameters after warm-up.
After one gradient step on new-domain data with learning rate $\eta$, the second-order approximation of old-task risk increase is:
\begin{equation}
    \mathcal{R}_{\mathrm{old}}(\theta^* + \delta) - \mathcal{R}_{\mathrm{old}}(\theta^*) \approx \underbrace{g_{\mathrm{old}}^\top \delta}_{\text{first-order}} + \tfrac{1}{2}\,\underbrace{\delta^\top H_{\mathrm{old}}\, \delta}_{\text{second-order}}
\label{eq:risk_increase}
\end{equation}
where $g_{\mathrm{old}} = \nabla_\theta \mathcal{R}_{\mathrm{old}}(\theta^*)$, $H_{\mathrm{old}} = \nabla^2_\theta \mathcal{R}_{\mathrm{old}}(\theta^*)$, and $\delta = -\eta\, g_{\mathrm{new}}$ is the parameter update.

\textbf{(a) Shared-parameter architecture.}
Partition $\theta = [\theta_{\mathrm{shared}};\, \theta_{\mathrm{spec}};\, \theta_{\mathrm{spa}}]$. The Hessian $H_{\mathrm{old}}$ contains cross-block terms $H_{\mathrm{shared,spa}}$ and $H_{\mathrm{shared,spec}}$ that couple all branches.
Adapting $\theta_{\mathrm{spa}}$ to the new domain propagates through $H_{\mathrm{shared,spa}}$ to damage the shared representation, increasing old-task risk for \emph{all} branches.

\textbf{(b) Independent-branch architecture with spectral freezing.}
When branches share no parameters, the feature Jacobians $J_b = \partial f_b / \partial \theta_b$ are block-diagonal, meaning the Hessian is \emph{approximately} block-diagonal:
\begin{equation}
    H_{\mathrm{old}} \approx \mathrm{diag}(H_{\mathrm{spec}},\, H_{\mathrm{hsi}},\, H_{\mathrm{lid}})
\end{equation}
The approximation is exact when the loss decomposes additively over branches.
In practice, the classifier operates on concatenated features $f = [\fspec;\, \fhsi;\, \flid]$, introducing residual cross-branch coupling through the shared loss surface. The magnitude of this coupling depends on the classifier's sensitivity to cross-branch feature interactions and is not analytically bounded here.
With the spectral branch frozen ($\delta_{\mathrm{spec}} = 0$), the old-task risk increase is dominated by spatial-branch terms:
\begin{equation}
    \mathcal{R}_{\mathrm{old}}(\theta^* + \delta) - \mathcal{R}_{\mathrm{old}}(\theta^*) \approx \sum_{b \in \{\mathrm{hsi, lid}\}} \left(g_{\mathrm{old},b}^\top \delta_b + \tfrac{1}{2}\,\delta_b^\top H_b\, \delta_b\right) + O(\epsilon_{\mathrm{cross}})
\label{eq:independent_risk}
\end{equation}
where $\epsilon_{\mathrm{cross}}$ captures the residual cross-branch coupling through the classifier.
The spectral branch's contribution is \emph{approximately zero}---because $\delta_{\mathrm{spec}} = 0$ and cross-branch terms are small for independent architectures.
\end{proposition}

\textbf{Interpretation and limitations.}
The key qualitative insight is robust: in independent architectures, adapting spatial parameters causes substantially less collateral damage to the spectral representation than in shared architectures.
The block-diagonal Hessian is an approximation---the shared classifier introduces residual coupling---but the empirical drift measurements in Section~\ref{sec:obs_results} confirm the predicted pattern: independent architectures show drift ratios of $1.5$--$5.8\times$ (spatial drift dominates, spectral is stable), while coupled architectures show ratios of $0.4$--$1.6\times$.

\subsubsection{Remark: Margin-Preserving Adaptation via Spectral Anchoring + DCR}
\label{app:margin}

We sketch an intuitive argument showing that spectral anchoring combined with DCR helps preserve the \emph{classification margin} of old classes under domain-selective LoRA adaptation. The following analysis uses a simplified single-centroid cosine scorer for tractability; the qualitative conclusion extends to the multi-centroid scorer used in practice, since each centroid component undergoes the same LoRA transformation.

\begin{remark}[Old-class margin preservation, simplified scorer]
\label{remark:margin}
Let $s(c \mid x) = \cosim(f(x),\, \mu_c)$ be the cosine similarity score for a single centroid (the bound applies to each centroid component in our multi-centroid scorer).
Define the margin of sample $x$ with true class $y$ as $m(x) = s(y \mid x) - \max_{c \neq y} s(c \mid x)$.
After adapting to a new domain with LoRA update $\Delta W_b$ on spatial branches ($b \in \{\mathrm{hsi, lid}\}$) and applying DCR, the margin perturbation satisfies:
\begin{equation}
    |m'(x) - m(x)| \leq \frac{4\,\epsilon_{\mathrm{spa}}}{(1 - \epsilon_{\mathrm{spa}})^2}
\label{eq:margin_bound}
\end{equation}
where $\epsilon_{\mathrm{spa}} = \max_{b \in \{\mathrm{hsi,lid}\}} \|\Delta W_b\|$ controls the spatial perturbation magnitude, and $m'$ is the margin under the adapted model with DCR-corrected prototypes.
\end{remark}

\begin{proof}[Sketch]
After DCR correction, both the query feature and the prototype undergo the same linear transformation $(I + \Delta W_b)$ on spatial branches, with the spectral component unchanged.
Write the full adapted feature as $\tilde{f} = [\fspec;\, (I+\Delta W_{\mathrm{hsi}})\fhsi;\, (I+\Delta W_{\mathrm{lid}})\flid]$ and similarly for prototypes.

The cosine similarity between adapted query and DCR-corrected prototype can be expressed as:
\begin{equation}
    \cosim(\tilde{f}, \tilde{\mu}_c) = \frac{\tilde{f}^\top \tilde{\mu}_c}{\|\tilde{f}\|\,\|\tilde{\mu}_c\|}
\end{equation}

Since the spectral component is identical in both numerator terms and $\Delta W_b$ has spectral norm $\epsilon_b \leq \epsilon_{\mathrm{spa}}$, we bound $\|\tilde{f} - f\| \leq \epsilon_{\mathrm{spa}} \|f_{\mathrm{spa}}\|$ where $f_{\mathrm{spa}} = [\fhsi;\,\flid]$.
Applying the Lipschitz property of cosine similarity (which has Lipschitz constant $\leq 2/\|f\|$ for unit-normed arguments), the per-class score perturbation satisfies $|s'(c|x) - s(c|x)| \leq \frac{2\epsilon_{\mathrm{spa}}}{1-\epsilon_{\mathrm{spa}}}$, and the margin perturbation is at most twice the per-class bound. \qed
\end{proof}

\textbf{Key insight.}
Without DCR, the query features pass through $(I + \Delta W_b)$ but old prototypes remain stale, creating a \emph{systematic} mismatch that grows with $\|\Delta W_b\|$ and is not captured by the symmetric bound above.
DCR eliminates this asymmetry, ensuring that the margin perturbation depends only on the \emph{second-order} effect of the transformation on cosine geometry, not on a first-order feature-prototype misalignment.

Without spectral freezing, $\epsilon_{\mathrm{spa}}$ would include spectral-branch updates, increasing the perturbation bound by the ratio of total to spatial-only update norms.
The combination of spectral anchoring (reduces the set of perturbed dimensions) and DCR (eliminates asymmetric mismatch) is what keeps $|m' - m|$ small.

\subsubsection{Proposition 3: Asymmetric Rank Allocation Under Budget Constraint}
\label{app:rank_allocation}

\begin{proposition}[Rank allocation under budget constraint]
\label{prop:rank}
For each spatial branch $b \in \{\mathrm{hsi, lid}\}$, let $\Delta W_b^* \in \R^{d \times d}$ be the ideal (full-rank) weight update that perfectly compensates the domain shift, with singular values $\sigma_1^{(b)} \geq \sigma_2^{(b)} \geq \cdots \geq \sigma_d^{(b)}$.
A rank-$r_b$ LoRA adapter incurs approximation error (by the Eckart--Young--Mirsky theorem):
\begin{equation}
    \mathcal{E}_b(r_b) = \|\Delta W_b^* - [\Delta W_b^*]_{r_b}\|_F^2 = \sum_{i=r_b+1}^{d} (\sigma_i^{(b)})^2
\label{eq:eckart_young}
\end{equation}
Given a total rank budget $R = r_{\mathrm{hsi}} + r_{\mathrm{lid}}$, the allocation minimizing $\mathcal{E}_{\mathrm{hsi}}(r_{\mathrm{hsi}}) + \mathcal{E}_{\mathrm{lid}}(r_{\mathrm{lid}})$ satisfies the KKT condition:
\begin{equation}
    (\sigma_{r_{\mathrm{hsi}}+1}^{(\mathrm{hsi})})^2 = (\sigma_{r_{\mathrm{lid}}+1}^{(\mathrm{lid})})^2
\label{eq:kkt_rank}
\end{equation}
That is, under the stated linearization assumptions, the \emph{marginal} tail energy should be equalized across branches at the budget-constrained optimum. In practice, training dynamics and nonlinearities introduce deviations; the proposition provides a principled initialization heuristic validated by the rank allocation comparison (Table~\ref{tab:rank_allocation}).
\end{proposition}

\begin{proof}
The objective is to minimize $\mathcal{E}(r_{\mathrm{hsi}}, r_{\mathrm{lid}}) = \sum_{i > r_{\mathrm{hsi}}} (\sigma_i^{(\mathrm{hsi})})^2 + \sum_{i > r_{\mathrm{lid}}} (\sigma_i^{(\mathrm{lid})})^2$ subject to $r_{\mathrm{hsi}} + r_{\mathrm{lid}} = R$ with $r_b \in \mathbb{Z}_{\geq 0}$.
Moving one unit of rank from $b_1$ to $b_2$ changes the objective by $(\sigma_{r_{b_1}}^{(b_1)})^2 - (\sigma_{r_{b_2}+1}^{(b_2)})^2$.
At optimum, no transfer improves the objective, yielding the equalized marginal condition~\eqref{eq:kkt_rank}. \qed
\end{proof}

\textbf{Connection to asymmetric rank allocation.}
Proposition~\ref{prop:rank} predicts that branches whose ideal update $\Delta W_b^*$ has a slower-decaying singular spectrum (energy spread across more directions) require higher rank.
Table~\ref{tab:rank_allocation} validates this empirically: at the same total budget ($R{=}12$), our asymmetric allocation ($r_{\mathrm{hsi}}{=}4$, $r_{\mathrm{lid}}{=}8$) outperforms the symmetric baseline ($6,6$) by $1.1$\,pp.

\textbf{Intuition for asymmetric allocation.}
The relevant quantity is \emph{effective adaptation dimensionality}, not CKA drift magnitude.
CKA measurements show comparable or even lower drift for LiDAR-spatial than HSI-spatial on certain domains (e.g., lid/hsi drift ratio ${\approx}0.5$ on MUUFL).
Yet the LiDAR branch needs more rank because it is the \emph{sole carrier} of elevation information with no spectral branch backup, requiring adaptation across more diverse directions even when the total change magnitude is moderate.
Table~\ref{tab:rank_allocation} validates this: at the same budget, our asymmetric allocation ($4,8$) outperforms the symmetric baseline ($6,6$).

\subsubsection{Remark: SHINE as Affine Nuisance Removal}
\label{app:shine_remark}

SHINE normalization (Eq.~\ref{eq:shine_norm}) can be understood as removing an affine domain nuisance.
Model the domain shift as an element-wise affine transformation $f_b^{(d)} = a_d \odot f_b + b_d$ where $a_d, b_d \in \R^d$ are domain-specific scale and bias.
Z-score normalization with domain statistics $(\mu_d^b, \sigma_d^b)$ inverts this nuisance exactly when $a_d = \sigma_d^b / \sigma_0^b$ and $b_d = \mu_d^b - a_d \odot \mu_0^b$ for some reference domain 0.
This holds per-coordinate and does not require Gaussianity.

The +5.5pp gain from SHINE (84.5\% $\to$ 79.0\% without) indicates that the affine component is the dominant cross-domain nuisance in this setting---consistent with the neural collapse literature showing that deep features exhibit approximately affine class-conditional structure~\cite{papyan2020prevalence}.


\subsection{Note on 9-Task Drift Consistency}
\label{app:9task_drift}

The drift observation study (Section~\ref{sec:observation}) uses a simplified 3-task diagnostic protocol (one task per domain) to isolate cross-domain drift from within-domain class-incremental drift.
Under the full 9-task protocol used in the main experiments, within-domain class-incremental steps are expected to introduce minor additional drift, but the cross-domain asymmetry pattern should remain consistent with the 3-task diagnostic, since the dominant drift source is the geographic domain shift rather than within-domain class addition.
A systematic 9-task drift analysis is left for future work.

\subsection{Confusion Matrices}
\label{app:confusion}

Confusion matrices reveal that \methodname{}'s errors are primarily within-domain misclassifications (e.g., confusing similar vegetation classes within MUUFL), rather than cross-domain confusion.
This confirms that the frozen spectral anchor and SHINE normalization effectively prevent cross-domain feature collapse, while the remaining errors reflect genuine spectral/spatial similarity between land-cover types within the same geographic region.


\bibliographystyle{IEEEtran}
\bibliography{references}   

\end{document}  
